\section{Three parameter certificates}\label{sec:certificates}

We now select three scripts in \(\calF_7\) and specify their leaf laws.
The first two establish the dual-exponent and rectangular conclusions
of \cref{thm:rectangular}.  The third records the square bound
\(\omega<2.267\) achieved by this finite family, independently of the
stronger square construction in \cref{thm:square}.
The inequalities in this section use rational enclosures, not numerical
feasibility assumptions.

\subsection{A dual-exponent witness}\label{subsec:dual}

Use base two, distinguished side \(X\), and its color \(\colB\).
After \(b\) base factors, regard the \(X\)-coordinates as binary words
of length \(b\); the logarithmic base cost is \(b\log2\).
We seek a final leaf law for which \(X\) is uniform on all these words
and \(S=\HH(X)=b\log2\).
By \cref{cor:balance} and the output-size lower bound, such a law gives
\[
 w\!\left(\min\{1,2T/(b\log2)\}\right)=2.
\]
The construction will preserve the exact incident-entropy equality
while increasing the nonincident entropy \(T\).

To identify what preserves this equality, let \(X\) also denote the
random coordinate induced by a leaf law on a readable tree, including
a tree conditional on one color.
At a node \(\nu\), conditioning means that the path reaches that node,
and the conditional mutual information is
\[
 \II(X;\text{child}\mid\nu)
 =\HH(X\mid\nu)-\HH(X\mid\text{child},\nu).
\]
Since a leaf determines \(X\), successive conditioning at the
nonterminal nodes telescopes to
\[
 \HH(X)=\sum_\nu\Prb(\nu)\,
                  \II(X;\text{child}\mid\nu).
\]
At an incident node the child is a function of \(X\) and the prefix,
so its mutual information is \(\HH(\text{child}\mid\nu)\).
Subtracting these incident contributions gives
\begin{equation}\label{eq:information-leakage}
 \HH(X)-S
  =\sum_{\nu\text{ nonincident}}
       \Prb(\nu)\,\II(X;\text{child}\mid\nu).
\end{equation}
The terms on the right are nonnegative.  Thus \(S=\HH(X)\) precisely
when every nonincident choice at a positive-mass node is independent
of \(X\), conditional on the preceding choices.

We maintain this equality separately within each color.
Partition the binary words into an active set \(P\) and its complement
\(Q\), with the following properties:
\begin{enumerate}
 \item color \(\colB\) permits precisely \(P\), while each other color
 permits precisely \(Q\);
 \item the chosen conditional law of each color makes \(X\) uniform on
 its permitted set;
 \item its conditional incident entropy is exactly the logarithm of
 that set's size.
\end{enumerate}
Write \(\rho=|P|/2^b\), and let \(h_i\) be color \(i\)'s nonincident
entropy; initially \((\rho,b,h_i)=(1/2,1,0)\).

\begin{lemma}\label{lem:uniform-update}
The maintained properties are preserved by the following laws.
At a completion centered on \(\colB\) of size \(m\),
\begin{align}
 \rho'&=\rho^m,& h'_{\colB}&=m h_{\colB}, \label{eq:minus-rho}\\
 h'_i&=
 m\left(\frac{\rho-\rho'}{1-\rho'}h_{\colB}
         +\frac{1-\rho}{1-\rho'}h_i\right)
       &&(i=\colA,\colC). \label{eq:minus-h}
\end{align}
At a completion centered on \(i\in\{\colA,\colC\}\), with \(j\) the
other inactive color,
\begin{align}
 \rho'&=1-(1-\rho)^m,&
 h'_{\colB}&=
 m\left(\frac{\rho}{\rho'}h_{\colB}
       +\frac{\rho'-\rho}{\rho'}h_i\right), \label{eq:plus-rho}\\
 h'_i&=mh_i,&
 h'_j&=\log\left[(e^{h_i}+e^{h_j})^m-e^{mh_i}\right].
 \label{eq:plus-h}
\end{align}
In both cases \(b'=mb\).
\end{lemma}

\begin{proof}
At a \(\colB\)-centered completion, the new active set is \(P^m\).
A uniform word in its complement determines the \(\colB/i\) pattern
from its \(P/Q\) membership in each slot.
This pattern is an incident label.
Conditioning on not all slots being active, the expected proportions
of old \(\colB\) and old \(i\) slots are
\[
 \frac{\rho-\rho^m}{1-\rho^m},
 \qquad \frac{1-\rho}{1-\rho^m}.
\]
Use independent old conditional laws in those slots.
They give the required uniform \(X\)-law and
\eqref{eq:minus-h}.
Their old nonincident choices retain independence of the relevant
\(X\)-coordinates, so \eqref{eq:information-leakage} gives the incident
entropy equality.

At a completion centered on inactive \(i\), the new inactive set is
\(Q^m\).  For new color \(\colB\), the \(\colB/i\) pattern is again
determined by \(X\); conditioning on at least one active slot gives
the expected proportion \(\rho/\rho'\) of old \(\colB\) slots.
This proves \eqref{eq:plus-rho}.

For new color \(j\), every slot has an old color in \(\{i,j\}\), with
at least one \(j\).
Both old colors have exactly the same uniform \(X\)-marginal on \(Q\).
Consequently every such pattern has the same uniform marginal on
\(Q^m\), independently of the law chosen for the pattern.
The pattern label is nonincident, and is independent of the entire
\(X\)-word.
Optimize its entropy plus the old nonincident scores by
\eqref{eq:gibbs}, with weight one.
This gives \eqref{eq:plus-h}.
The old independent slot laws preserve the zero terms on the right
of \eqref{eq:information-leakage}.
Pure-center words use only the old \(i\)-law, proving the remaining
formula.
\end{proof}

Terminate by keeping \(\colB\) and the inactive color with larger \(h\),
placing mass \(\rho\) on \(\colB\).
The incident termination label and the uniform conditional laws make
\(X\) uniform on all \(2^b\) binary words.
Hence
\begin{equation}\label{eq:flat-rates}
 S=b\log2,\qquad
 T=bD,\qquad
 D=\frac{\rho h_{\colB}+(1-\rho)\max(h_{\colA},h_{\colC})}{b}.
\end{equation}
By \cref{cor:balance} and the size lower bound,
\begin{equation}\label{eq:alpha-certificate}
 \alpha\ge \min\{1,2D/\log2\}.
\end{equation}

Use the following seven completions:
\begin{equation}\label{eq:dual-script}
 -2,\quad +2,\quad +2,\quad -3,\quad +2,\quad -2,\quad +3.
\end{equation}
A minus sign means center \(\colB\).
A plus sign means the inactive color with larger \(h\), with an
arbitrary fixed tie choice.
The absolute value is the completion size.
Resolving the first tie as \(\colA\), the centers are
\(\colB,\colA,\colC,\colB,\colC,\colB,\colA\).

For a short certificate, set
\(\delta=|h_{\colA}-h_{\colC}|\).
A minus step leaves \(D\) unchanged and gives
\[
 \delta'=\frac{m(1-\rho)}{1-\rho^m}\delta.
\]
At a plus step put
\[
 z=\log\bigl((1+e^{-\delta})^m-1\bigr).
\]
Substitution in \cref{lem:uniform-update} yields
\begin{equation}\label{eq:delta-D}
 \delta'=|z|,\qquad
 D'=D+\frac{1-\rho'}{mb}\max\{0,z\}.
\end{equation}
Indeed \(h_i\) is the larger inactive entropy, and
\(h'_j=mh_i+z\), so the new maximum is
\(mh_i+\max\{0,z\}\).

After the first three steps, at \(b=8\),
\[
 \rho_8=1-(3/4)^4,\qquad
 \delta_8=\log(9/7),\qquad
 D_8=\frac{0.5625}{4}\log3.
\]
The two plus steps here have \(z=\log3\) and \(z=\log(7/9)\).
The subsequent bounds are
\[
\begin{array}{c|l}
 b&\text{certified information}\\ \hline
24&
 \rho_{24}=\rho_8^3,\quad
 0.3194<\rho_{24}<0.3195,\quad
 0.348<\delta_{24}<0.354\\
48&
 1-\rho_{48}=(1-\rho_{24})^2>0.463,\quad
 0.536<\rho_{48}<0.537,\quad 0.635<z_{48}<0.651\\
96&
 \rho_{96}=\rho_{48}^2,\quad
 \delta_{96}=2\delta_{48}/(1+\rho_{48})<0.86\\
288&
 1-\rho_{288}=(1-\rho_{48}^2)^3>0.360,\quad z_{288}>0.628.
\end{array}
\]
Their elementary verification is given in \cref{subsec:log-certificate}.
Since \(\log2<0.6932\) and \(\log3>1.098\),
\begin{align}
 \frac{2D_{288}}{\log2}
 &>
 \frac{2}{0.6932}
 \left[
  \frac{0.5625\cdot1.098}{4}
  +\frac{0.463\cdot0.635}{48}
  +\frac{0.360\cdot0.628}{288}
 \right]\notag\\
 &=\frac{1548637}{3327360}>0.465. \label{eq:alpha-rational}
\end{align}
This proves the dual-exponent conclusion without any approximate
equality being used to certify exponent two.

\subsection{A one-parameter rectangular curve}\label{subsec:curve}

Use base three, one size-three completion centered on \(\colA\),
keep \(\colB,\colC\), and distinguish \(X\).
This is the 75-component support of \cref{thm:elementary}; the leaf
law below tunes its two balanced outer dimensions and its inner dimension.
Every word in \(\{0,1,2\}^3\) occurs on \(X\).
For a nonzero word \(x\), its outer color is \(\colB\), and the number
of supported leaves in its fiber is
\(d(x)=2^{\#\{\text{zero entries of }x\}}\).
For the zero word, its outer color is \(\colC\), and
\(d(0)=3^3-2^3=19\).
The fiber histogram is therefore
\[
\begin{array}{c|rrrr}
 d&1&2&4&19\\ \hline
 \#\{x:d(x)=d\}&8&12&6&1.
\end{array}
\]
The total number of leaves is \(75\).

Give \(x\) probability proportional to \(d(x)^t\), and choose a leaf
uniformly within its fiber.
Put
\begin{equation}\label{eq:curve-functions}
 \begin{split}
 Z(t)&=8+12\,2^t+6\,4^t+19^t,\\
 g(t)&=(\log Z)'(t),\qquad h(t)=\log Z(t)-tg(t).
 \end{split}
\end{equation}
Then \(\HH(X)=h(t)\) and \(\HH(\text{leaf}\mid X)=g(t)\).
In the nonzero branch, the \(\colA/\colB\) pattern and the choices
within \(\colB\) are incident to \(X\); the choices within \(\colA\)
are hidden on \(Y,Z\).
In the zero branch all choices after branching are on \(Y,Z\).
The outer \(\colB/\colC\) label is incident to \(X\).
To check that the hidden choices do not leak information about \(X\),
fix a nonzero outer branch and its \(\colA/\colB\) pattern.
Then \(d(x)=2^{\#\colA}\) is constant over that pattern.
Each leaf has probability \(d(x)^{t-1}/Z(t)\), so the hidden choices
in its \(\colA\)-slots are independent of the remaining
\(X\)-coordinates, including after any preceding slot histories.
The zero branch has fixed \(X\).  The information identity
\eqref{eq:information-leakage} therefore gives precisely
\(S=h(t)\), \(T=g(t)\).
Its base cost is \(3^3\), and \cref{cor:balance} gives
\begin{equation}\label{eq:curve-bound}
 w\!\left(\frac{2g(t)}{h(t)}\right)
       \le\frac{6\log3}{h(t)} .
\end{equation}
Here \(w(k)=W(1,k,1)\) can also be read for \(k>1\);
we use only aspect ratios in \([0,1]\).

At \(t=2/3\), write \(u=2^{2/3}\), \(v=19^{2/3}\).
Then
\[
 g=\frac{12u(1+u)\log2+v\log19}{8+12u+6u^2+v}.
\]
Cubing verifies
\[
 1.58739<u<1.58742,\qquad
 7.1203<v<7.1205,\qquad
 49.2878<Z<49.2890.
\]
Using the logarithm bounds below gives
\[
 1.11844<g<1.11856,\qquad 3.1519<h<3.1522.
\]
In particular
\[
 \frac{2g}{h}>
 \frac{2\cdot1.11844}{3.1522}>0.709,\qquad
 \frac{6\log3}{h}<
 \frac{6\cdot1.0987}{3.1519}<2.092 .
\]
Monotonicity proves \(w(0.709)<2.092\).

\subsection{A finite-family square witness}\label{subsec:square}

Use base three and the completions
\[
 (\colB,2),\quad(\colA,3),\quad(\colC,2),\quad
 (\colB,2),\quad(\colA,3),
\]
then keep \(\colB,\colC\).
The number of original factors is \(b=72\).
Let \(d_i\) be the number of conditional monomials of color \(i\).
Initially \((d_{\colB},d_{\colA},d_{\colC})=(2,2,1)\).
A size-\(m\) completion at \(i\) updates counts to
\[
 d_i^m\quad\text{at }i,\qquad
 (d_i+d_j)^m-d_i^m\quad\text{at }j\ne i.
\]
After three steps the vector is
\((20691584,13993344,10144225)\).
After four it is
\[
 (428141648429056,\;774902581936128,\;522705468255425)
 \ \ge\ 10^{12}(428,774,522)
\]
coordinatewise.
The final sum of the kept counts is
\[
 N=2995461856366867388223311649574852466276480577,
\]
and in particular
\begin{equation}\label{eq:square-count-lower}
 N\ge 10^{36}
      (1202^3+1296^3-2\cdot774^3)>2.95\cdot10^{45}.
\end{equation}
The substitution of lower counts is legitimate: the polynomial
\((a+b)^3+(a+c)^3-2a^3\) has nonnegative coefficients.

Under the uniform leaf law, the sum of the three pair rates is
\(\log N\).
Averaging over permutations balances all three at \(\log N/3\).
Thus \cref{thm:entropy} gives
\[
 w(1)\le \frac{3\cdot72\log3}{\log N}
 <\frac{216\cdot1.0987}{1.0817+45\cdot2.3025}
 <2.267.
\]

\noindent The three constructions above yield their stated bounds
once the following elementary enclosures are verified.

\subsection{Elementary interval verification}\label{subsec:log-certificate}

For \(1\le y\le2\), put \(s=(y-1)/(y+1)\).
The logarithm series has the rational enclosure
\begin{equation}\label{eq:log-series}
 2\sum_{j=0}^7\frac{s^{2j+1}}{2j+1}
 \le\log y\le
 2\sum_{j=0}^7\frac{s^{2j+1}}{2j+1}
       +\frac{2s^{17}}{17(1-s^2)}.
\end{equation}
It follows by summing the positive series for
\(2\operatorname{arctanh}s\) and bounding its tail geometrically.
For a positive rational \(x\), write \(x=2^a y\) with \(1\le y\le2\)
and use \(\log x=a\log2+\log y\).
Rational substitution in \eqref{eq:log-series} gives
\[
\begin{array}{c|c}
x&\text{bounds for }\log x\\ \hline
2&(0.69314,\;0.69316)\\
3&(1.098,\;1.0987)\\
19&(2.9444,\;2.9445)\\
{[49.2878,49.2890]}&(3.89765,\;3.89772)\\
9/7&(0.250,\;0.252)\\
10&>2.3025\\
2.95&>1.0817.
\end{array}
\]
For the dual witness,
\[
 \delta_{24}=
 \frac{3(1-\rho_8)}{1-\rho_8^3}\log(9/7);
\]
the displayed bounds for \(\rho_{24}\) and \(\delta_{24}\) now use
only rational arithmetic.
Alternating Taylor bounds through degrees six and seven give
\[
 e^{-0.354}>0.701,\qquad
 e^{-0.348}<0.707,\qquad
 e^{-0.86}>0.4225.
\]
Consequently the logarithm series at
\[
 (1+0.701)^2-1,\quad
 (1+0.707)^2-1,\quad
 (1+0.4225)^3-1
\]
proves \(0.635<z_{48}<0.651\) and \(z_{288}>0.628\).
Also
\[
 \delta_{96}<\frac{2\cdot0.651}{1+0.536}<0.86.
\]
The remaining \(\rho\) inequalities are direct rational powers.
This verifies every inequality used in \eqref{eq:alpha-rational}.
The cube-root brackets and the expression for \(g\) similarly verify
the curve certificate.

For reference only, evaluation of the exact formulas gives
\[
\begin{array}{c|c}
\text{quantity}&\text{approximate value}\\ \hline
2D_{288}/\log2&0.466101554898\\
2g(2/3)/h(2/3)&0.709698932519\\
6\log3/h(2/3)&2.091251835621\\
216\log3/\log N&2.266187438999.
\end{array}
\]
The proofs above rely on the rational inequalities, not these decimal
approximations.
The accompanying certificate scripts reproduce the finite
support/count calculations and all stated rational enclosures.
In a fresh writable copy of \texttt{verification/certificates/}, the command
\texttt{python3 checks/run\_checks.py}
uses integer and rational arithmetic from the Python standard library;
its three checkers separately verify the numerical inequalities,
finite supports and counts, and declared baseline data.
The optional numerical diagnostics are not used in these certificates.
